%% P34 --- Uczciwy pomiar modelu.  OSTATNI PROGRAM W KSIĄŻCE.
%%
%% Bliźniak: programs/en/P34-measuring-honestly.tex.  Nagłówek angielski
%% niesie cały odczyt briefu i listę programów, które odsyłają tutaj z
%% nazwy (F01, F04, F08, P17) --- nie powtarzam go, bo kopia to następna
%% rzecz, która się zdezaktualizuje.
%%
%% TŁUMACZENIE, decyzje zapisane tutaj, żeby nikt ich nie relitygował:
%%   benchmark        -> benchmark, odmieniany po polsku (reguła CLAUDE.md
%%                       o angielskich terminach ML bez polskiej formy)
%%   judge model      -> model sędziujący, dalej krótko sędzia
%%   attenuation      -> tłumienie
%%   agreement        -> zgodność
%%   probe            -> sonda, za Programem P31, który tę decyzję zapisał
%%   prompt           -> prompt, odmieniany
%%   leaderboard      -> tablica wyników
%%
%% UWAGA DLA TŁUMACZA: C4, C8 i C12 porównują znaczniki strukturalne, wzory i
%% literały liczbowe PO KOLEI. Cyfra zostaje cyfrą, a słowo zostaje słowem;
%% odwołanie do programu prowadzi zdanie zamiast stać za wzorem.

\program{Uczciwy pomiar modelu}
\label{prog:P34}
\index{ewaluacja}

\begin{outcomes}
\outcome{1--7}{powiedzieć, czego oszacowaniem jest wynik benchmarku i którą
  połowę jego błędu da się zmniejszyć liczbą pozycji}
\outcome{8--15}{przeważyć jeden zbiór ewaluacyjny i powiedzieć, co to robi z
  uporządkowaniem, którego nikt nie uważał za sporne}
\outcome{16--19}{powiedzieć, ile pozycji kosztuje zwężenie raportowanego
  przedziału}
\outcome{20--28}{powiedzieć, co niedoskonały sędzia robi ze zmierzoną różnicą
  i dlaczego jego zgodność tego nie zdradza}
\outcome{29--35}{wycenić model za zadanie, a nie za token, i powiedzieć,
  której drugiej wielkości brakuje w cenie za token}
\outcome{36--43}{wymienić trzy publikowane miary policzone poprawnie z
  niewłaściwej wielkości}
\outcome{44--47}{powiedzieć, które tezy tego programu nie zostały sprawdzone
  na wytrenowanym modelu i dlaczego powiedzenie tego nie jest asekuracją}
\outcome{48--52}{wymienić trzy powody, dla których tabela w książce jest
  pusta, i powiedzieć, który z nich jest długiem}
\end{outcomes}

\begin{quiz}
\item Twój zbiór ewaluacyjny ma $\val{p27.items}$ pozycji, a model dostaje na
  nim $\val{p27.near.hi}$ procent. O czym jest ta liczba?
  \teachesat{1--2}
  \answerto{To oszacowanie skuteczności na populacji zadań, na której ci
  zależy. Zbiór jest próbą, a wynik estymatorem, więc obie osie z
  Programu~\ref{prog:P26} i przedział z Programu~\ref{prog:P27} stosują się
  do niego bez zmian.}
\item Którą połowę błędu benchmarku zmniejsza większa liczba pozycji?
  \teachesat{3--5}
  \answerto{Wyłącznie połowę losową. Dystans między tym, co mierzy benchmark,
  a zadaniem, na którym ci zależy, nie maleje z liczbą pozycji wcale, a
  powyżej około $\val{p34.mis.n}$ pozycji jest większy z tych dwóch.}
\item Dwa modele dzieli $\val{p27.gap.pts}$ punktu na jednym zbiorze
  ewaluacyjnym. Oba zostają bez zmian, a proporcja krótkich i długich promptów
  w zbiorze zmienia się o $\val{p34.mix.shift}$ pozycji. Co się dzieje?
  \teachesat{11--14}
  \answerto{Uporządkowanie się odwraca. Wyniki stają się
  $\val{p34.mix.alt.b}$ i $\val{p34.mix.alt.a}$, a własny wynik jednego
  modelu przebiega $\val{p34.mix.swing}$ punktów po wszystkich mieszankach.}
\item Model sędziujący ocenia odpowiedzi niedoskonale. Co to robi z różnicą,
  którą raportuje między dwoma modelami?
  \teachesat{21--23}
  \answerto{Mnoży ją przez jedną stałą, tę samą dla każdej pary. Jeśli sędzia
  uznaje dobrą odpowiedź za dobrą w $\val{p34.judge.a}$ procentach przypadków,
  a błędną w $\val{p34.judge.b}$, każda zmierzona różnica wynosi
  $\val{p34.judge.atten}$ prawdziwej.}
\item Model B kosztuje $\val{p34.bill.price.times}$ raza tyle za token co
  model A. Który jest droższy?
  \teachesat{30--31}
  \answerto{Z tego nie da się odpowiedzieć. Cena za token to połowa ilorazu:
  rachunek to cena razy liczba tokenów, które pochłania zadanie, a przy
  $\val{p34.bill.tok.b}$ tokenach wobec $\val{p34.bill.tok.a}$ model B jest o
  $\val{p34.bill.saving.pct}$ procent tańszy na zadanie.}
\item Sonda raportuje $\val{p34.probe.i}$ nata informacji wzajemnej między
  warstwą a etykietą, na $50$ pozycjach. Ile z tego jest estymatorem?
  \teachesat{38--39}
  \answerto{$\val{p34.probe.share.pct}$ procent.
  Program~\ref{prog:P31} zmierzył, że estymator wtykowy zwraca
  $\val{p31.bias.n50}$ nata, w wartości oczekiwanej, na danych bez żadnej
  zależności.}
\end{quiz}

% ============================================================
\section{Zbiór ewaluacyjny jest próbą}
\label{sec:P34-a-sample}
\index{ewaluacja!jako zagadnienie estymacji}

\begin{fr}
Program~\ref{prog:P27} zostawił ci tablicę wyników: $\val{p27.items}$ pozycji,
najlepszy model na $\val{p27.near.hi}$ procentach, następny na
$\val{p27.near.lo}$, i dokładny test mówiący, że różnica między nimi jest
najmniej informatywnym wynikiem, na jaki arytmetyka pozwala.

Ten program jest o wszystkim innym na tamtej stronie. Zacznij od samego
wyniku. Uruchomiłeś ustalony zbiór pozycji, policzyłeś i podzieliłeś. Liczba
wynosi $\val{p27.near.hi}$ i jest dokładną arytmetyką na tym, co uruchomiłeś.

Czego więc dokładnie jest \emph{oszacowaniem}?
\dotline
\nextframe
\end{fr}

\begin{fr}
\ans{Skuteczności na populacji zadań, na której naprawdę ci zależy.}

Nikt nie raportuje wyniku benchmarku dlatego, że zależy mu na tych dwustu
pozycjach. Raportuje go dlatego, że pozycje mają zastąpić pewną pracę. To
czyni ze zbioru \textbf{próbę}, a z wyniku \textbf{estymator}, i wprowadza do
pokoju cały Program~\ref{prog:P26}.

Program~\ref{prog:P26} rozdzielił błąd estymatora na dwie części i dowiódł
tego rozdziału dokładnie. Nazwij je.
\dotline
\nextframe
\end{fr}

\begin{fr}
\ans{Kwadrat obciążenia plus wariancja.}

Odczytane na zbiorze ewaluacyjnym obie części mają nazwy, których inżynier
już używa.

\textbf{Wariancja} to losowanie pozycji: uruchom inne dwieście, a dostaniesz
inną liczbę. Program~\ref{prog:P27} posiada ją w całości i to właśnie ona jest
opisana przez przedział w raporcie.

\textbf{Obciążenie} to dystans między tym, co mierzy benchmark, a pracą, do
której zatrudniasz model. Jego pozycje napisał ktoś inny, dla cudzego ruchu.

Jedna z tych dwóch maleje, gdy dokładasz pozycji. Która?
\dotline
\nextframe
\end{fr}

\begin{fr}
\ans{Wyłącznie wariancja.}

Wariancja proporcji wynosi $p(1-p)/n$, więc maleje jak $1/n$ i da się ją
zbić tak nisko, jak pozwoli budżet. Obciążenie jest własnością \emph{zbioru},
nie jego rozmiaru: dziesięć tysięcy pozycji napisanych z tego samego
błędnego założenia jest błędne dokładnie tak samo jak dwieście.

Weź model o prawdziwej skuteczności $\num{0.715}$, na benchmarku odległym o
$\val{p34.mis.pts}$ punkty od twojego zadania. Przy jakiej liczbie pozycji
obie części wnoszą tyle samo do błędu tego, co raportujesz?
\dotline
\nextframe
\end{fr}

\begin{fr}
\begin{ansblock}
Przyrównaj obie i rozwiąż względem $n$:
\[ \frac{p(1-p)}{n} = \text{obciążenie}^{2}
   \qquad\Longrightarrow\qquad
   n = \frac{p(1-p)}{\text{obciążenie}^{2}}
     = \frac{\num{0.715} \times \num{0.285}}{\num{0.02}^{2}} , \]
co przechyla się po raz pierwszy przy $n = \val{p34.mis.n}$.
\end{ansblock}

To nie jest duża liczba. To około dwa i pół raza tyle, co zbiór, o który
spierał się Program~\ref{prog:P27}.

Program~\ref{prog:P27} uruchomił własną tablicę na $\val{p27.lb.n}$
pozycjach. Która z dwóch części wnosi tam więcej i o ile?
\dotline
\nextframe
\end{fr}

\begin{fr}
\ans{Niedopasowanie, mniej więcej $\val{p34.mis.ratio}$ raza.}

Przy tysiącu pozycji dystans benchmarku od twojego zadania wnosi z grubsza
dwa razy tyle, co losowanie \dash{} i \textbf{tylko połowa losowa pojawia się
w przedziale w raporcie}. Druga połowa nie jest mała, nie jest raportowana i
nie rusza się, gdy dokupujesz pozycje.

Stąd zdanie, dla którego istnieje ta sekcja: \textbf{więcej pozycji kupuje
precyzję i nigdy nie kupuje trafności.}

Zostaje pytanie, które przekazał tu z nazwy Program~\ref{prog:F08}. Skoro
przedział nie powie ci, czy wynik cokolwiek znaczy na twoich danych, co
powie?
\dotline
\nextframe
\end{fr}

\begin{fr}
\begin{ansblock}
Wyłącznie porównanie pozycji benchmarku z twoim własnym ruchem. To praca
czytelnicza i praca losująca, nie arytmetyczna: pobierz próbę z tego, co
twój system naprawdę dostaje, i zapytaj, czy pozycje benchmarku wyglądają
jak próba z tego samego źródła.
\end{ansblock}

Warto powiedzieć to wprost, bo to jedyna część ewaluacji, której nie da się
kupić. Pod każdym innym pytaniem w tym programie leży wzór. Pod tym leży
decyzja projektowa, a decyzja projektowa, której nikt nie zapisał, jest
obciążeniem, którego nikt nie oszacuje.

\mermaidfig{p34-two-errors}{Dwa błędy raportowanego wyniku i ten, który
  opisuje przedział}{dwa błędy, jeden przedział}
\end{fr}

% ============================================================
\section{Te same pozycje, inaczej ważone}
\label{sec:P34-remix}
\index{ewaluacja!mieszanka jako wybór}

\begin{fr}
Teraz zostaw wszystko w obu modelach bez zmian i nie ruszaj niczego w tym, co
odpowiedziały. Poruszy się wyłącznie \emph{zbiór}.

Program~\ref{prog:P27} dostarcza $\val{p27.items}$ pozycji; podziel je na dwa
rodzaje: $\val{p34.mix.short}$ krótkich promptów i $\val{p34.mix.short}$
długich. Każdy benchmark ma podział tego kształtu, niezależnie od tego, czy
go raportuje.

Model B trafia $\val{p34.mix.b.short}$ procent krótkich pozycji i
$\val{p34.mix.b.long}$ procent długich. Model A trafia
$\val{p34.mix.a.short}$ i $\val{p34.mix.a.long}$ \dash{} lepiej na długich
promptach, gorzej na krótkich. Policz oba na całym zbiorze.
\dotline
\nextframe
\end{fr}

\begin{fr}
\begin{ansblock}
Połowa pozycji jest każdego rodzaju, więc każdy wynik to zwykła średnia:
\[ B = \tfrac{1}{2}(\num{0.80}) + \tfrac{1}{2}(\num{0.63}) = \val{p27.near.hi}\%,
   \qquad
   A = \tfrac{1}{2}(\num{0.74}) + \tfrac{1}{2}(\num{0.68}) = \val{p27.near.lo}\% . \]
\end{ansblock}

To są dwie zapisane liczby z Programu~\ref{prog:P27}, co do cyfry. To ta sama
tablica wyników, z jedną rzeczą o niej zapisaną, która była tam od początku.

Równy podział był wyborem, którego ktoś dokonał. Załóż, że zamiast tego
przeważyliby krótkie prompty wagą $w$. Co dzieje się z różnicą?
\dotline
\nextframe
\end{fr}

\begin{fr}
\ans{Porusza się po prostej, bo wynik to średnia ważona.}

Wynik każdego modelu jest liniowy w $w$, więc różnica między nimi też:
\[ B - A = (\num{0.63} - \num{0.68}) + w\bigl[(\num{0.80}-\num{0.74}) -
   (\num{0.63}-\num{0.68})\bigr] = -\num{0.05} + \num{0.11}\,w . \]

Przy $w = \frac{1}{2}$ daje to $+\num{0.005}$, czyli tę różnicę, którą
Program~\ref{prog:P27} podaje jako $\val{p27.gap.pts}$ punktu.

Gdzie się zeruje?
\dotline
\nextframe
\end{fr}

\begin{fr}
\ans{Przy $w = \frac{5}{11}$, czyli $\val{p34.mix.tie.w}$ procentach krótkich
  promptów.}

To nie jest dziwna mieszanka. To zbiór z odrobinę większą liczbą długich
promptów niż krótkich, czyli to, co wyszłoby z pisania pozycji innego
popołudnia.

Równy podział daje $\val{p34.mix.short}$ krótkich pozycji w zbiorze. Ile z
$\val{p27.items}$ musi zmienić rodzaj, zanim czołowa dwójka tablicy zamieni
się miejscami?
\dotline
\nextframe
\end{fr}

\begin{fr}
\begin{ansblock}
Dziesięć. Przy $\val{p34.mix.shift}$ krótkich pozycji mniej wyniki wynoszą
\[ B = \val{p34.mix.alt.b}\% , \qquad A = \val{p34.mix.alt.a}\% , \]
więc A prowadzi teraz o $\val{p34.mix.alt.gap}$ punktu.
\end{ansblock}

Żaden z modeli nie został dotrenowany, przestrojony ani inaczej
promptowany. Każda odpowiedź, której którykolwiek z nich udzielił, jest ta
sama co przedtem. Dziesięć pozycji na dwieście przeszło z jednego kubełka do
drugiego i uporządkowanie, które raportuje tablica, jest odwrotne.

\transcript{p34-remix}
\end{fr}

\begin{fr}
Różnica jest tu małą wielkością, warto więc zapytać, jak duża jest ta
\emph{druga}.

Sam model B, po wszystkich mieszankach od samych długich do samych krótkich:
jaki jest zakres wyników, które można mu przypisać?
\dotline
\nextframe
\end{fr}

\begin{fr}
\ans{Od $\val{p34.mix.b.long}$ do $\val{p34.mix.b.short}$ procent \dash{}
  $\val{p34.mix.swing}$ punktów.}

Postaw to obok dwóch liczb, o które tablica naprawdę się spiera. Omawiana
różnica to $\val{p27.gap.pts}$ punktu. Program~\ref{prog:P27} policzył, że ten
zbiór ewaluacyjny nie odróżni obu modeli, dopóki różnica nie sięgnie
$\val{p27.net.needed.pts}$ punktu.

\begin{trapbox}
Wynik jednego modelu porusza się po mieszankach $\val{p34.mix.swing.times}$
raza dalej niż różnica, która według Programu~\ref{prog:P27} byłaby cokolwiek
warta.

Mieszanka nie jest hiperparametrem, nie ma jej w przedziale, a na większości
tablic wyników nie jest raportowana. Wielkość, która rusza wynikiem bardziej
niż jakakolwiek różnica między modelami i której nikt nie zapisuje, to
definicja zmiennej niekontrolowanej.
\end{trapbox}
\end{fr}

\begin{fr}
Nic z tego nie jest zarzutem wobec benchmarku. Zbiór pozycji \emph{musi} mieć
jakąś mieszankę; neutralnej nie ma, a odmowa wyboru nie jest dostępna.

Znaczy to coś węższego i jest to zdanie z Programu~\ref{prog:F01}, które
właśnie przychodzi. Wynik nie jest \enquote{tym, jak dobry jest model}. Jest
\emph{tym, jak dobry na tej mieszance}, a mieszanka jest tą drugą wielkością.
Podaj pierwszą bez drugiej, a podałeś iloraz bez jednej z jego dwóch liczb
\dash{} czyli kształt, który ta książka nazywa od swojego pierwszego programu.

\mermaidfig{p34-same-items}{Jeden zbiór ewaluacyjny, dwie mieszanki i oba
  modele nietknięte}{średnia ważona rusza się z wagami}
\end{fr}

% ============================================================
\section{Ile kosztuje węższy przedział}
\label{sec:P34-interval}
\index{ewaluacja!koszt precyzji}

\begin{fr}
Program~\ref{prog:P27} zapisał jeden przedział, a ta sekcja go wycenia. Pozycja
tablicy na $\val{p27.lb.p.pct}$ procentach przy $\val{p27.lb.n}$ pozycjach ma
błąd standardowy $\val{p27.lb.se}$ punktu, a przy tamtejszej poprawce na
czytanie czterdziestu pozycji margines wynosi $\val{p27.lb.margin}$ punktu.

Załóż, że to za szeroko, żeby działać, i chcesz połowy. Ile pozycji?
\dotline
\nextframe
\end{fr}

\begin{fr}
\ans{$\val{p34.halve.times}$ raza tyle \dash{} $\val{p34.halve.n}$.}

Precyzja idzie jak $1/\sqrt{n}$, więc połowienie przedziału to podniesienie
czynnika do kwadratu. Warto powiedzieć to w walucie, w której się za to
naprawdę płaci: cztery razy tyle pozycji do napisania, cztery razy tyle do
oznaczenia i czterokrotny rachunek za wnioskowanie, które je oceni.

Pociągnij o krok dalej. Ile pozycji, żeby podać skuteczność jednego modelu z
dokładnością do punktu, przy tej samej poprawce?
\dotline
\nextframe
\end{fr}

\begin{fr}
\begin{ansblock}
Liczba skaluje się jak kwadrat marginesu, który kupuje:
\[ n = \val{p27.lb.n} \times
   \left(\frac{\val{p27.lb.margin}}{\num{1.0}}\right)^{2}
   \approx \val{p34.one.pt.n} . \]
\end{ansblock}

Blisko dziesięciu tysięcy pozycji, żeby podać jedną liczbę z dokładnością do
punktu.

Program~\ref{prog:P27} policzył z tej samej maszynerii inną liczbę: ile
pozycji potrzeba, żeby \emph{odróżnić dwa modele} o punkt, gdy oba biegną na
tych samych promptach. Jak się te dwie mają do siebie?
\dotline
\nextframe
\end{fr}

\begin{fr}
\ans{$\val{p27.n.rho00}$ wobec $\val{p34.one.pt.n}$ \dash{} czynnik
  $\val{p34.detect.ratio}$.}

Podanie skuteczności jednego modelu z dokładnością do punktu i oddzielenie
dwóch modeli o punkt kosztują mniej więcej tyle samo, a oba są dwa rzędy
wielkości za zbiorem, który masz przed sobą.

\begin{note}
To jest uczciwy odczyt przedziału: jest \textbf{deklaracją budżetu}. Mówi,
co jest w stanie rozstrzygnąć ewaluacja, za którą już zapłaciłeś, a
odpowiedź brzmi zwykle \enquote{mniej niż omawiana różnica}. Nikt nie ma
dziesięciu tysięcy pozycji na model, dlatego niemal każde publikowane
porównanie jest raczej niedomocowane niż błędne.
\end{note}
\end{fr}

% ============================================================
\section{Sędzia i to, co robi z różnicą}
\label{sec:P34-judge}
\index{ewaluacja!model sędziujący}

\begin{fr}
Dziesięć tysięcy pozycji to nie jest budżet na oznaczanie, który ktokolwiek
ma, więc pozycje ocenia inny model.

Sędzia jest przyrządem, a przyrząd ma dwie stopy błędu, nie jedną. Zapisz
\[ a = \Prob(\text{sędzia uznaje za dobrą} \mid \text{naprawdę dobra}),
   \qquad
   b = \Prob(\text{sędzia uznaje za dobrą} \mid \text{naprawdę błędna}). \]

Prawdziwa skuteczność modelu to $p$. Jaką skuteczność raportuje dla niego
sędzia?
\dotline
\nextframe
\end{fr}

\begin{fr}
\ans{$p\,a + (1-p)\,b$.}

Sędzia uznaje za dobrą ułamek $a$ dobrych odpowiedzi i błędnie uznaje za
dobrą ułamek $b$ błędnych. Wnoszą obie.

Nikt jednak nie raportuje jednej skuteczności samej. Dwa modele o prawdziwych
skutecznościach $p_{1}$ i $p_{2}$ przechodzą przez tego samego sędziego. Co
raportuje dla \emph{różnicy} między nimi?
\dotline
\nextframe
\end{fr}

\begin{fr}
\begin{ansblock}
Odejmij jedno od drugiego, a wszystko, co nie niesie $p$, się skraca:
\[ \bigl[p_{1}a + (1-p_{1})b\bigr] - \bigl[p_{2}a + (1-p_{2})b\bigr]
   = (a - b)(p_{1} - p_{2}) . \]
\end{ansblock}

To jest dokładne i warto zatrzymać się na dwóch rzeczach.

Raportowana różnica to prawdziwa różnica pomnożona przez \textbf{jedną
stałą}, więc niedoskonały sędzia nie dodaje szumu do siły efektu \dash{} on
ją \emph{kurczy}, systematycznie, w jedną stronę. A stała $a - b$ \textbf{nie
zależy od $p$ wcale}, więc to samo kurczenie stosuje się do porównania na
szczycie tablicy i do porównania na jej dole.
\end{fr}

\begin{fr}
Podstaw liczby. Sędzia uznaje dobrą odpowiedź za dobrą w
$\val{p34.judge.a}$ procentach przypadków, a błędną odpowiedź za dobrą w
$\val{p34.judge.b}$, więc $a - b = \val{p34.judge.atten}$. Dwa modele dzieli
naprawdę $\val{p34.judge.true.gap}$ punktu, a sędzia raportuje
$\val{p34.judge.seen.gap}$.

Prawdziwej różnicy nie zobaczysz, musisz więc wykryć raportowaną. Co kurczenie
robi z liczbą pozycji, której potrzebujesz?
\dotline
\nextframe
\end{fr}

\begin{fr}
\ans{Mnoży ją przez $\val{p34.judge.items.times}$.}

Efekt kurczy się o $a - b$, a szum losowy nie, więc iloraz obu kurczy się
również o $a - b$; a arytmetyka z Sekcji~\ref{sec:P34-interval} mówi, że
liczba skaluje się jak kwadrat rozdzielczości, którą kupuje. Liczba pozycji
dzieli się więc przez $(a-b)^{2}$ \dash{} o połowę więcej pozycji z okładem,
przy dokładnie tej samej mocy, kupione niczym innym niż ocenianiem.

To nie jest mała poprawka do budżetu. Nie ma jej też w niczyim raporcie.
\end{fr}

\begin{fr}
Chciałbyś więc znać $a - b$ dla sędziego, którego używasz. To, co podaje
karta sędziego, to \textbf{zgodność} z etykietami człowieka: ułamek pozycji,
na których sędzia i człowiek mówią to samo, co dla modelu o skuteczności $p$
wynosi
\[ \text{zgodność} = p\,a + (1-p)(1-b) . \]

Przy skuteczności z tego programu i tych dwóch stopach daje to
$\val{p34.judge.agree.pct}$ procent, czyli przyzwoitą kartę.

Czy znajomość zgodności zdradza tłumienie?
\dotline
\nextframe
\end{fr}

\begin{fr}
\begin{ansblock}
Nie. Zgodność wiąże kombinację $p\,a - (1-p)\,b$ i nic więcej, więc przy
$\val{p34.judge.agree.pct}$ procentach z kartą zgodny jest każdy sędzia o
\[ \val{p34.judge.atten.lo} \;\le\; a - b \;\le\;
   \val{p34.judge.atten.hi} . \]
\end{ansblock}

Oba końce różnią się na tyle, żeby to miało znaczenie: przez $(a-b)^{2}$
dzieli je $\val{p34.judge.spread.pct}$ procent w liczbie pozycji, której
żądają. Dwaj sędziowie o \emph{identycznej} zgodności na tej samej
oznaczonej próbie, a między nimi połowa budżetu na oznaczanie z okładem.
\end{fr}

\begin{fr}
\begin{trapbox}
Istnieje dokładnie jedna skuteczność, przy której zgodność wiąże tłumienie,
i warto wiedzieć która.

Obie stopy wchodzą do zgodności ze współczynnikami $p$ i $-(1-p)$, więc
równoważą się \dash{} a kombinacja zwija się do $a - b$ \dash{} tylko przy
$p = \frac{1}{2}$. Wtedy zgodność wyznacza tłumienie dokładnie, jako
$2\times\text{zgodność} - 1$.

Jedyny przypadek, w którym karta sędziego wystarcza, to więc przypadek, w
którym model trafia połowę czasu, czyli jedyna skuteczność, o której nikt nie
raportuje tablicy wyników.
\end{trapbox}
\end{fr}

\begin{fr}
Program~\ref{prog:P28} też wycenił sędziego i warto być precyzyjnym co do
różnicy, bo te dwie poprawki nie zastępują się nawzajem.

Tamten program bierze prawdopodobieństwo, które sędzia \emph{podaje} \dash{}
mówi $\val{p28.judge.says.pct}$ procent, a znaczy $\val{p28.judge.means.pct}$
\dash{} i poprawia je w szansach. Ta sekcja bierze \emph{werdykt} sędziego i
poprawia to, co robi on ze zmierzoną różnicą. Jedno jest stwierdzeniem o
liczbie, którą sędzia podaje o sobie; drugie stwierdzeniem o wpływie przyrządu
na twój eksperyment.

Skalibrowanie pierwszego zostawia drugie dokładnie tam, gdzie było.

\mermaidfig{p34-judge-shrinks}{Niedoskonały sędzia mnoży każdą siłę efektu
  przez jedną stałą}{werdykt, nie podane prawdopodobieństwo}
\end{fr}

% ============================================================
\section{Rachunek, który jest iloczynem}
\label{sec:P34-bill}
\index{koszt!na zadanie wobec na token}

\begin{fr}
Program~\ref{prog:F01} otworzył tę książkę stwierdzeniem, że niemal każda
teza, którą przeczytasz o modelu \dash{} szybszy, tańszy, lepszy \dash{} jest
ilorazem podanym bez jego dwóch wielkości, i odłożył czytanie takich tez
tutaj.

Koszt za token ma dokładnie ten kształt. To jedna wielkość, a to, za co
ktokolwiek naprawdę płaci, jest iloczynem dwóch.

Jaka jest ta druga?
\dotline
\nextframe
\end{fr}

\begin{fr}
\ans{Liczba tokenów, które pochłania zadanie.}

\[ \text{koszt na zadanie} \;=\; \text{cena za token} \times
   \text{tokeny na zadanie} . \]

Nikt nie kupuje tokenów. Kupuje odpowiedziane pytania, zamknięte zgłoszenia,
przejrzane diffy \dash{} a model, który mówi więcej, żeby dojść do tej samej
odpowiedzi, wydaje różnicę.

Model B kosztuje $\val{p34.bill.price.times}$ raza tyle za token co model A i
zużywa średnio $\val{p34.bill.tok.b}$ tokenów na zadanie wobec
$\val{p34.bill.tok.a}$ u A. Który jest droższy w utrzymaniu?
\dotline
\nextframe
\end{fr}

\begin{fr}
\begin{ansblock}
Pomnóż oba ilorazy:
\[ \frac{\text{B na zadanie}}{\text{A na zadanie}}
   = \val{p34.bill.price.times} \times \val{p34.bill.tok.times}
   = \val{p34.bill.task.times} , \]
więc B jest o $\val{p34.bill.saving.pct}$ procent \emph{tańszy} na zadanie,
będąc o połowę droższym za token.
\end{ansblock}

Oba ilorazy ciągną w przeciwne strony i rozliczany jest wyłącznie ich
iloczyn.

Przy jakiej rozwlekłości odpowiedź by się odwróciła \dash{} ile tokenów na
zadanie mógłby B sobie pozwolić, zanim tańszy stałby się A?
\dotline
\nextframe
\end{fr}

\begin{fr}
\ans{$\val{p34.bill.breakeven}$ tokenów.}

B musi być co najmniej o trzecią bardziej zwięzły od A, żeby spłacić swoją
cenę, i przy $\val{p34.bill.tok.b}$ tokenach jest. Przesuń go do
$\val{p34.bill.breakeven}$, a oba stają dokładnie równo; przesuń dalej, a
liczba za token znowu ma rację.

Żadna z liczb nie rozstrzyga więc niczego sama, a próg opłacalności jest
własnością pary, nie któregokolwiek z modeli.
\end{fr}

\begin{fr}
Za tymi samymi dwiema wielkościami stoi trzecia i ma ją
Program~\ref{prog:P03}. Obsługa rozmowy trzyma bufor
$\val{p03.kv.per.token.mib}$ MiB na token w locie, więc długość zadania
kosztuje pamięć tak samo jak pieniądze \dash{} $\val{p34.bill.cache.a}$ MiB
dla A i $\val{p34.bill.cache.b}$ dla B.

Cena za token nie pojawia się w tym wcale. O czym to decyduje?
\dotline
\nextframe
\end{fr}

\begin{fr}
\ans{O tym, ile rozmów mieści się naraz na jednej karcie \dash{} B obsługuje
  $\val{p34.bill.conc.times}$ raza tyle.}

Bufor na token jest ten sam dla obu modeli, więc skraca się z ilorazu w
całości, a zostaje długość zadania. Zwięźlejszy model jest tańszy na zadanie
\emph{i} niesie więcej równoczesnych użytkowników na tym samym sprzęcie, a
karta mówiąca, że jest $\val{p34.bill.price.times}$ raza droższy, nie wspomina
o żadnym z tych dwóch.
\end{fr}

\begin{fr}
Program~\ref{prog:P29} spotkał tę tożsamość w innej walucie i warto postawić
obie obok siebie, bo zachowują się inaczej.

Tam bity na znak to bity na token razy tokeny na znak i \textbf{wkład modelu
się skraca}: ten sam model na tym samym tekście raportuje $\val{p29.bpc.a}$ i
$\val{p29.bpc.b}$ bita na znak przy $\val{p29.tpc.a}$ i $\val{p29.tpc.b}$
tokena na znak, czyli czynnik $\val{p29.bpc.ratio}$, który jest faktem o
tokenizatorze i o niczym innym.

Tutaj nie skraca się żaden z czynników. I cena, i długość należą do modelu,
dlatego właśnie podanie jednej z nich niczego nie załatwia.

\mermaidfig{p34-two-quantities}{Cena za token to połowa tego, co kosztuje
  zadanie}{rozliczany jest tylko iloczyn}
\end{fr}

% ============================================================
\section{Miary, które wyglądają na pomiary}
\label{sec:P34-measures}
\index{ewaluacja!miary, które mylą}

\begin{fr}
Trzy publikowane wielkości są policzone poprawnie i odpowiadają na pytanie,
którego nikt nie zadał. Dwie z nich ta książka już zmierzyła i warto zebrać
je tutaj, bo kształt, który dzielą, jest sednem.

Program~\ref{prog:P29} wziął straty dwóch modeli przeliczone na bity na znak.
Przy różnych tokenizatorach co skraca się z tego porównania?
\dotline
\nextframe
\end{fr}

\begin{fr}
\ans{Model.}

Bity na znak to bity na token razy tokeny na znak, a drugi czynnik należy do
tokenizatora. Porównaj dwa systemy, które tokenizują inaczej, a różnica,
którą zmierzysz, należy do tokenizatorów. Dwie straty są porównywalne przy
jednym tokenizatorze i pod żadnym innym warunkiem, a podanie tego warunku
kosztuje zdanie.
\end{fr}

\begin{fr}
Ta z Programu~\ref{prog:P31} jest ostrzejsza, bo liczba nie jest tam jedynie
niejednoznaczna \dash{} jest dodatnia, kiedy prawda wynosi zero.

Sonda raportuje $\val{p34.probe.i}$ nata informacji wzajemnej między
aktywacjami warstwy a etykietą, na tablicy zbudowanej z $50$ pozycji. Zanim
cokolwiek z tego odczytasz: co ten estymator raportuje na danych bez żadnej
zależności?
\dotline
\nextframe
\end{fr}

\begin{fr}
\begin{ansblock}
$\val{p31.bias.n50}$ nata \dash{} w wartości oczekiwanej, dokładnie, przy tej
liczbie pozycji. Program~\ref{prog:P31} zsumował to po każdej tablicy,
zamiast losować.
\end{ansblock}

Zatem $\val{p34.probe.share.pct}$ procent raportowanej liczby to estymator, a
nie warstwa. Efekt może być prawdziwy; jedna piąta z niego nie jest.

Obciążenie maleje jak $1/N$. Ile pozycji, zanim spadnie poniżej jednej
dwudziestej tamtej raportowanej liczby?
\dotline
\nextframe
\end{fr}

\begin{fr}
\ans{Około $\val{p34.probe.n}$.}

Cztery razy tyle pozycji, żeby zbić artefakt do jednej dwudziestej tezy, na
pomiarze, który rutynowo raportuje się z kilkudziesięciu. I w
przeciwieństwie do błędu losowego ten nie okrąża prawdy \dash{} idzie w
jedną stronę, więc nigdy nie uśrednia się po warstwach jednego artykułu.
\end{fr}

\begin{fr}
Program~\ref{prog:P31} niesie też trzecią i jest to jedna nierówność czytana
na dwa końce.

Sonda wypada słabo na warstwie. Czy to znaczy, że warstwa niesie mało o
etykiecie?
\dotline
\nextframe
\end{fr}

\begin{fr}
\ans{Nie. Wynik sondy to ograniczenie dolne, i to luźne.}

Program~\ref{prog:P31} zmierzył najlepszy kanał w rodzinie sięgający
$\val{p31.dpi.best}$ nata wobec $\val{p31.dpi.ixy}$, które tam były, czyli z
ponad ćwiercią zostawioną. Ta sama nierówność czytana w drugą stronę mówi, że
\emph{lepsza} sonda nie może oznaczać więcej informacji, bo sufit był ustalony,
zanim ktokolwiek jakąkolwiek wytrenował.

Jedna liczba jest więc rutynowo nadinterpretowana w obie strony, przez tych
samych ludzi, w tym samym artykule.
\end{fr}

\begin{fr}
Trzy miary, trzy różne sposoby na to, by nie być pomiarem tego, do czego się
je przywołuje, i jeden kształt między nimi: \textbf{wielkość policzona
poprawnie z niewłaściwej rzeczy.}

To są cztery ciche awarie z Programu~\ref{prog:P33} przybywające o poziom
wyżej. Tam krok treningowy mógł policzyć średnią ze średnich albo zaokrągloną
sumę i nie podnieść niczego, bo średnia ze średnich jest średnią. Tutaj
benchmark może raportować bity na znak, albo estymatę wtykową sondy, albo
skuteczność na niezapisanej mieszance, a każde z nich jest arytmetyką, której
nikt nie zarzuci błędu.

Błędna odpowiedź sama się ogłasza. Poprawna odpowiedź na niewłaściwe pytanie
nie, i to właśnie czyni z projektu całą robotę.
\end{fr}

% ============================================================
\section{Czego ta książka nie sprawdziła}
\label{sec:P34-not-checked}
\index{pomiar!czego się nie twierdzi}

\begin{fr}
Program~\ref{prog:P17} odłożył tu z nazwy jedną tezę: czy płaskie minimum
generalizuje lepiej niż ostre. To najczęściej powtarzana teza empiryczna w
optymalizacji, a ta książka nie powiedziała o niej nic.

Zadaj jej pytanie z Programu~\ref{prog:P32}. Która połowa rozstrzygnięcia
wymaga wytrenowanego modelu, a która jest stwierdzeniem o projekcie?
\dotline
\nextframe
\end{fr}

\begin{fr}
\begin{ansblock}
Całość wymaga wytrenowanych modeli i to jest nietypowa odpowiedź.

Teza porównuje dwa \emph{wytrenowane} minima, więc nie ma czego wyprowadzać:
potrzeba wielu modeli, definicji płaskości, która przeżyje przeskalowanie z
Programu~\ref{prog:P17}, i tylu przebiegów na warunek, ile zażądałby
Program~\ref{prog:P27}.
\end{ansblock}

Jest więc podana jako zaległa, a nie zasekurowana.
Programy~\ref{prog:P08}, \ref{prog:P11}, \ref{prog:P19}, \ref{prog:P20},
\ref{prog:P25}, \ref{prog:P31} i \ref{prog:P32} niosą po jednej takiej i
każdy mówi, co by ją rozstrzygnęło.
\end{fr}

\begin{fr}
To samo pytanie ten program jest winien o swoich własnych trzech wynikach.
Który z przeważenia, tłumienia sędziego i rachunku wymaga wytrenowanego
modelu?
\dotline
\nextframe
\end{fr}

\begin{fr}
\ans{Żaden.}

Każdy jest stwierdzeniem o \emph{projekcie}, a nie o modelu. Przeważenie to
arytmetyka na średniej ważonej; tłumienie to tożsamość na dwóch
prawdopodobieństwach warunkowych; rachunek to iloczyn dwóch ilorazów. Każde
jest dokładne i żadne nie byłoby prawdziwsze z modelem za sobą.

\begin{warning}
Wytrenowanego modelu wymagałaby połowa empiryczna: czy prawdziwe zbiory
ewaluacyjne mają mieszanki tak wpływowe jak ta skonstruowana w
Sekcji~\ref{sec:P34-remix} i jak daleko od siebie leżą tłumienia prawdziwych
sędziów przy równej zgodności.

Obie są mierzalne i żadnej tu nie zmierzono. Pierwszą rozstrzygnąłby
publikowany benchmark z pozycjami skategoryzowanymi i przeważonymi; druga
wymaga oznaczonej próby ocenionej przez dwóch sędziów.
\end{warning}
\end{fr}

% ============================================================
\section{Standard}
\label{sec:P34-standard}
\index{pomiar!standard, który liczba musi spełnić}

\begin{fr}
Oba tomy towarzyszące tej książce niosą tabelę, w której nic nie ma, i ta
książka również. Jest to zamierzone we wszystkich trzech i nie jest jedną
decyzją powtórzoną.

Zanim przeczytasz dalej: tabela w książce technicznej jest pusta. Wymień
powody, dla których może być, i powiedz, który z nich jest długiem, który ktoś
jest ci winien.
\dotline
\nextframe
\end{fr}

\begin{fr}
\begin{ansblock}
Trzy, i tylko pierwszy jest długiem.

\begin{enumerate}
\item \textbf{Eksperymentu nie przeprowadzono.} Dług \dash{} a książka jest ci
  winna to, co by go rozstrzygnęło, i ile to kosztuje.
\item \textbf{Liczba jest własnością twojej maszyny.} Nie dług. Tabela
  wypełniona przez kogoś innego byłaby czynnie myląca, a to, co tam należy,
  to metoda.
\item \textbf{Zmierzono ją i oznaczono jako orientacyjną}, bo rusza się z
  wyborem, którego nikt nie ustandaryzował.
\end{enumerate}
\end{ansblock}

Tomy towarzyszące niosą po jednym z każdego. Tabela dostawców w tomie o Agent
Framework jest tym drugim: jej własna nota mówi, że liczby zależą od twojego
sprzętu i twojego regionu, więc metodologia jest ustalona, a liczby są twoje.
Tom o LangChain drukuje pierwszy i trzeci na jednej stronie \dash{}
wypełnioną tabelę, której kolumna kontekstu jest zatytułowana
\emph{orientacyjna}, obok trzech pustych wierszy mówiących, że wypełnienie ich
prawdopodobnymi liczbami byłoby gorsze niż zostawienie ich pustymi.
\end{fr}

\begin{fr}
Warto obchodzić się z tymi dwiema książkami ostrożnie, zamiast je zlewać, bo
obie drukują liczbę $\num{2.8}$ raza, a te liczby nie są tym samym pomiarem.

Jedna to wielokrotność kontekstu zmierzona na skryptowanych odpowiedziach i
oznaczona jako orientacyjna dokładnie z tego powodu. Druga to stosunek
tokenów, którego druga książka \emph{odmawia podania} bez zmierzenia,
mówiąc, że książka twierdząca to bez pomiaru byłaby gorsza niż taka, która
przyznaje, że nie zmierzyła.

Żadna nie traktuje tej liczby jako ustalonej, i to jest jedyne, co obie mają
wspólnego. Dwie książki, jedna liczba, dwa zupełnie różne powody, by jej nie
ufać \dash{} a odczytanie ich jako zgody byłoby popełnieniem na własnych
źródłach defektu z Sekcji~\ref{sec:P34-measures} tego programu.
\end{fr}

\begin{fr}
Oto więc standard, na którym zbudowane są oba tomy towarzyszące, wreszcie z
matematyką tej książki za sobą.

\begin{center}
\emph{Teza potrzebuje metody i liczby, albo jest oznaczona jako osąd.}
\end{center}

Za każde słowo z tego już zapłacono. \textbf{Liczba} jest oszacowaniem, więc
Program~\ref{prog:P26} mówi, jakie są jej dwa błędy, a
Program~\ref{prog:P27}, jak jest szeroka. \textbf{Metoda} jest projektem
\dash{} która populacja, która mieszanka, który przyrząd, który mianownik
\dash{} a ten program pokazał, że każde z tych czterech rusza raportowaną
liczbą dalej niż różnice, o które ludzie się spierają. A \textbf{osąd} jest
uczciwą etykietą dla tezy bez obu, co Program~\ref{prog:F04} miał na myśli,
czyniąc z tego standard, który raportowana liczba musi spełnić.
\end{fr}

\begin{fr}
To koniec matematyki.

Zacząłeś przy Programie~\ref{prog:F01}, wyliczając, ile waży model, bez
słownictwa na to, gdzie są jego parametry. Program~\ref{prog:P32} wyprowadził
tę liczbę z kształtów; Program~\ref{prog:P33} odczytał krzywą, którą rysuje
jego trening; a ten program mówi, co wolno o wyniku twierdzić.

Książka przez cały czas dowodziła, że matematyka nie jest w tej pracy częścią
trudną i nie jest też opcjonalna \dash{} że różnica między inżynierem, który
umie sprawdzić tezę, a takim, który ją powtarza, to kilkaset stron dokładnie
tego. Masz teraz te strony.

Zostaje nawyk, a jest nim jedno pytanie zadawane każdej liczbie, którą
ktokolwiek ci podaje, własnej także: \textbf{co zostało zmierzone i skąd
wiedziałbym, gdyby było błędne?}
\end{fr}

\begin{summarybox}
\begin{sumlist}
\sumitem{2--3}{Zbiór ewaluacyjny jest próbą, a wynik estymatorem, więc oba
  błędy z Programu~\ref{prog:P26} stosują się do niego: losowanie pozycji i
  dystans między benchmarkiem a twoim zadaniem.}
\sumitem{4--6}{Tylko pierwszy maleje z liczbą pozycji. Powyżej około
  $\val{p34.mis.n}$ pozycji drugi jest większy i nie ma go w przedziale:
  więcej pozycji kupuje precyzję i nigdy nie kupuje trafności.}
\sumitem{11--12}{Przeważenie jednego ustalonego zbioru ewaluacyjnego o
  $\val{p34.mix.shift}$ pozycji z $\val{p27.items}$ odwraca uporządkowanie
  jego czołowej dwójki, bez tknięcia któregokolwiek z modeli.}
\sumitem{13--14}{Wynik jednego modelu przebiega $\val{p34.mix.swing}$ punktów
  po możliwych mieszankach \dash{} $\val{p34.mix.swing.times}$ raza więcej niż
  różnica, która według Programu~\ref{prog:P27} byłaby cokolwiek warta.}
\sumitem{17--18}{Precyzja idzie jak $1/\sqrt{n}$, więc połowienie przedziału
  to $\val{p34.halve.times}$ raza tyle pozycji, a podanie jednej skuteczności
  z dokładnością do punktu około $\val{p34.one.pt.n}$ z nich.}
\sumitem{21--22}{Niedoskonały sędzia raportuje $(a-b)$ raza prawdziwą różnicę
  między dowolnymi dwoma modelami, dokładnie, a stała nie zależy od
  skuteczności żadnego z nich.}
\sumitem{24--26}{Dzieli to użyteczną liczbę pozycji przez $(a-b)^{2}$, a
  zgodność sędziego nie wyznacza $a-b$: przy
  $\val{p34.judge.agree.pct}$ procentach biegnie ono od
  $\val{p34.judge.atten.lo}$ do $\val{p34.judge.atten.hi}$.}
\sumitem{27--28}{Zgodność wiąże tłumienie tylko przy skuteczności połowy, a
  poprawka z Programu~\ref{prog:P28} do podanego przez sędziego
  prawdopodobieństwa zostawia tłumienie nietknięte.}
\sumitem{30--31}{Koszt na zadanie to cena za token razy tokeny na zadanie,
  więc model $\val{p34.bill.price.times}$ raza droższy za token może być o
  $\val{p34.bill.saving.pct}$ procent tańszy w utrzymaniu.}
\sumitem{33--34}{Bufor jest na token w locie, więc zwięźlejszy model obsługuje
  też $\val{p34.bill.conc.times}$ raza tyle rozmów na jednej karcie.}
\sumitem{37--39}{Bity na znak mierzą tokenizator, gdy tokenizatory się
  różnią, a wtykowa informacja wzajemna jest dodatnia na danych niezależnych
  \dash{} $\val{p31.bias.n50}$ nata przy $50$ pozycjach.}
\sumitem{41--42}{Wynik sondy to luźne ograniczenie dolne, więc słaba sonda nie
  jest dowodem małej informacji, a lepsza nie jest dowodem większej.}
\sumitem{44--47}{Teza o ostrości z Programu~\ref{prog:P17} wymaga
  wytrenowanych modeli i jest podana jako zaległa; własne trzy wyniki tego
  programu nie wymagają żadnego, bo każdy jest stwierdzeniem o projekcie.}
\sumitem{48--49}{Pusta tabela ma trzy powody i tylko jeden jest długiem: nie
  przeprowadzono, nieprzenoszalna, albo zmierzona i oznaczona jako
  orientacyjna.}
\sumitem{51--52}{Teza potrzebuje metody i liczby, albo jest oznaczona jako
  osąd \dash{} a metoda jest projektem tak samo jak arytmetyka.}
\end{sumlist}
\end{summarybox}

\begin{testexercises}
\item Benchmark jest odległy o $3$ punkty od twojego zadania, a twój model
  dostaje na nim $\num{0.60}$. Przy jakiej liczbie pozycji niedopasowanie
  wnosi tyle samo co losowanie?
  \answerto{Około $267$. Punkt przecięcia to $p(1-p)/\text{obciążenie}^{2} =
  \num{0.60} \times \num{0.40} / \num{0.03}^{2}$, a większe obciążenie
  przesuwa go w dół: im gorzej benchmark pasuje do twojego zadania, tym
  wcześniej więcej pozycji przestaje pomagać.}
\item Dwa modele dostają $\num{0.70}$ i $\num{0.66}$ na krótkiej połowie
  zbioru oraz $\num{0.62}$ i $\num{0.70}$ na długiej. Przy jakiej mieszance
  remisują?
  \answerto{Przy $\frac{2}{3}$ krótkich. Różnica wynosi
  $(\num{0.62}-\num{0.70}) + w[(\num{0.70}-\num{0.66}) -
  (\num{0.62}-\num{0.70})] = -\num{0.08} + \num{0.12}\,w$, co zeruje się przy
  $w = \frac{2}{3}$.}
\item Twój przedział to $\pm 4$ punkty przy $500$ pozycjach. Ile pozycji na
  $\pm 1$?
  \answerto{$8000$. Liczba skaluje się jak kwadrat stosunku marginesów, więc
  $500 \times 4^{2}$ \dash{} szesnaście razy tyle pozycji na czterokrotną
  rozdzielczość.}
\item Sędzia uznaje dobrą odpowiedź za dobrą w $\num{0.90}$ przypadków, a
  błędną w $\num{0.20}$. Dwa modele dzieli naprawdę $6$ punktów. Co jest
  raportowane i o jaki czynnik rośnie twój budżet na pozycje?
  \answerto{$\num{4.2}$ punktu, a budżet rośnie o
  $1/\num{0.70}^{2} \approx \num{2.04}$ \dash{} tłumienie wynosi
  $a - b = \num{0.70}$, a liczba idzie jak jego kwadrat.}
\item Model C ma $\num{0.80}$ ceny modelu D za token i zużywa $\num{1.30}$
  raza tyle tokenów na zadanie. Który jest tańszy na zadanie?
  \answerto{D, o około $4$ procent: $\num{0.80} \times \num{1.30} =
  \num{1.04}$. Tańszy token przegrywa z dłuższą odpowiedzią, dlatego żadna z
  tych liczb nie rozstrzyga sama.}
\item Tabela w książce jest pusta, a książka mówi, że liczba zależy od twojego
  sprzętu. Czy to dług?
  \answerto{Nie. To drugi z trzech powodów: liczby, która jest własnością
  maszyny czytelnika, nie da się dostarczyć, a tabela wypełniona przez kogoś
  innego byłaby czynnie myląca. To, co jest tam winne, to metoda, a nie
  liczba.}
\end{testexercises}

\canyou

\begin{furtherproblems}
\item Pokaż, że jeśli obciążenie benchmarku wynosi $\beta$, a jego pozycje
  pochodzą z ustalonej populacji, to żadna liczba pozycji nie zbije błędu
  średniokwadratowego raportowanej liczby poniżej $\beta^{2}$.
  \answerto{Błąd wynosi $\beta^{2} + p(1-p)/n$, a drugi składnik jest ściśle
  dodatni dla każdego skończonego $n$, więc całość leży ściśle powyżej
  $\beta^{2}$ i dąży do niego z góry. Podłogę ustala zbiór, a nie jego
  rozmiar.}
\item Zbiór ma $k$ rodzajów pozycji zamiast dwóch. Pokaż, że wynik nadal jest
  liniowy w wagach, i powiedz, co zastępuje pojedynczy punkt remisu.
  \answerto{Wynik to $\sum_{i} w_{i} s_{i}$ przy $\sum_{i} w_{i} = 1$, więc
  różnica to $\sum_{i} w_{i}(s_{i} - t_{i})$ \dash{} liniowa w $w$. Zbiór
  remisów to przecięcie tej hiperpłaszczyzny z sympleksem, więc przy więcej
  niż dwóch rodzajach remisy są całą ścianą, a nie punktem, i mieszanek
  odwracających uporządkowanie jest na ogół wiele.}
\item Dwaj sędziowie mają tłumienia $a_{1}-b_{1}$ i $a_{2}-b_{2}$, a jeden
  ocenia wyjście drugiego. Jakie jest tłumienie pary?
  \answerto{Iloczyn. Każdy etap mnoży ocalały efekt przez własną stałą, więc
  łańcuchowanie sędziów kumuluje kurczenie geometrycznie \dash{} dlatego
  zautomatyzowany potok oceniający jednym modelem i filtrujący drugim
  potrzebuje więcej pozycji, niż sugeruje którykolwiek etap.}
\item Model E ma $r$ raza cenę modelu F za token. Pokaż, że zbiór stosunków
  rozwlekłości, przy których E jest tańszy na zadanie, jest przedziałem, i
  powiedz, jaki jest jego koniec.
  \answerto{E jest tańszy, gdy $r \times (\text{tokeny}_{E}/\text{tokeny}_{F})
  < 1$, więc stosunek musi leżeć w $\intco{0}{1/r}$. Koniec to dokładnie
  $1/r$: model dwa razy droższy musi być co najmniej dwa razy zwięźlejszy i
  żadna jakość tej arytmetyki nie zmienia.}
\item Obciążenie wtykowe z Programu~\ref{prog:P31} maleje jak $1/N$, a
  raportowany efekt nie maleje wcale. Pokaż, że liczba pozycji, przy której
  obciążenie jest ustalonym ułamkiem $f$ efektu $I$, jest odwrotnie
  proporcjonalna do obu.
  \answerto{Obciążenie wynosi około $c/N$ dla stałej $c$ ustalonej przez
  kształt tablicy, więc $c/N \le fI$ daje $N \ge c/(fI)$. Połowienie efektu,
  który zgłaszasz, albo połowienie ułamka, który tolerujesz, każde podwaja
  liczbę pozycji \dash{} im mniejszy wynik, tym drożej go ustalić.}
\end{furtherproblems}
